% concatenated from camera_ready.tex
%% LaTeX Template for ISIT 2025
%%
%% by Stefan M. Moser, October 2017
%% (with minor modifications by Tobias Koch, November 2023 and Michèle Wigger, November 2024)
%% 
%% derived from bare_conf.tex, V1.4a, 2014/09/17, by Michael Shell
%% for use with IEEEtran.cls version 1.8b or later
%%
%% Support sites for IEEEtran.cls:
%%
%% http://www.michaelshell.org/tex/ieeetran/
%% http://moser-isi.ethz.ch/manuals.html#eqlatex
%% http://www.ctan.org/tex-archive/macros/latex/contrib/IEEEtran/
%%

\documentclass[conference,letterpaper]{IEEEtran}
%% depending on your installation, you may wish to adjust the top margin:
\addtolength{\topmargin}{9mm}

%%%%%%
%% Packages:
%% Some useful packages (and compatibility issues with the IEEE format)
%% are pointed out at the very end of this template source file (they are 
%% taken verbatim out of bare_conf.tex by Michael Shell).
%
% *** Do not adjust lengths that control margins, column widths, etc. ***
% *** Do not use packages that alter fonts (such as pslatex).         ***
%
\usepackage[cmex10]{amsmath} % Use the [cmex10] option to ensure complicance
                             % with IEEE Xplore (see bare_conf.tex)
\usepackage{graphicx} % Required for inserting images
\usepackage{amsthm}
\usepackage{amssymb}
\usepackage{amsfonts}
\usepackage{amsmath}
\usepackage{amsthm}
\usepackage{amssymb}
\usepackage{dsfont}

\DeclareMathOperator*{\argmin}{argmin}
\DeclareMathOperator*{\argmax}{argmax}

\newtheorem{definition}{Definition}
\newtheorem{lemma}{Lemma}
\newtheorem{theorem}{Theorem}
\newtheorem{corollary}{Corollary}
\newtheorem{proposition}{Proposition}
\newtheorem{problem}{Problem}
\newtheorem{conjecture}{Conjecture}
\newtheorem{remark}{Remark}
\newtheorem{assumption}{Assumption}

\newcommand{\rayleigh}[2]{\exp\left(-\left(\frac{#1}{r_0}\right)^{#2}\right)}
\newcommand{\tti}{\rightarrow\infty}
\newcommand{\graph}{\mathcal{G}}
\newcommand{\bigO}[1]{\mathcal{O}\left(#1\right)}
\newcommand{\preal}{\mathbb{R}_{>0}}
\newcommand{\real}{\mathbb{R}}
\newcommand{\prob}[1]{\mathbb{P}\left(#1\right)}
\newcommand{\red}[1]{\textcolor{red}{#1}}
\newcommand{\green}[1]{\textcolor{green}{#1}}
\newcommand{\erf}[0]{\text{erf}}
\newcommand{\ttz}[0]{\rightarrow 0}
\newcommand{\tto}[0]{\rightarrow}
\newcommand{\plim}[0]{p\mbox{-}\lim}
\newcommand{\plimsup}[0]{p\mbox{-}\limsup}
\newcommand{\pliminf}[0]{p\mbox{-}\liminf}
\newcommand{\Var}{\text{Var}}

%% Please note that the amsthm package must not be loaded with
%% IEEEtran.cls because IEEEtran provides its own versions of
%% theorems. Also note that IEEEXplore does not accepts submissions
%% with hyperlinks, i.e., hyperref cannot be used.

\interdisplaylinepenalty=2500 % As explained in bare_conf.tex


%%%%%%
% correct bad hyphenation here
\hyphenation{op-tical net-works semi-conduc-tor}

% ------------------------------------------------------------
\begin{document}
\title{Entropy and Distributed Source Coding of Connected Soft Random Geometric Graphs} 

% %%% Single author, or several authors with same affiliation:
% \author{%
%  \IEEEauthorblockN{Author 1 and Author 2}
% \IEEEauthorblockA{Department of Statistics and Data Science\\
%                    University 1\\
 %                   City 1\\
  %                  Email: author1@university1.edu}% }


%%% Several authors with up to three affiliations:
\author{%
  \IEEEauthorblockN{Oliver Baker and Carl P. Dettmann}
  \IEEEauthorblockA{School of Mathematics \\
                    University of Bristol\\
                    Bristol, United Kingdom\\
                    Email: \{o.baker, carl.dettmann\}@bristol.ac.uk}
}




\maketitle

\begin{abstract}
    We consider the distributed compression of Soft Random Geometric Graphs (SRGGs) above the connectivity threshold. We establish the Slepian-Wolf rate region for the SRGG in the setting where there are a finite number of encoders compressing sections of the graph independently. To do so, we prove novel limit theorems and asymptotic equipartition properties for the SRGG and its entropy, which allow us to use random binning techniques for distributed compression.
\end{abstract}

\section{Introduction}

In the modern world, devices are more connected than ever. This results in networks becoming larger, increasing the need for research on compression of networks. In particular, in areas of application such as wireless networking, random models of large networks with a spatial embedding have gained popularity \cite{haenggi_2012}. One popular model for wireless networks is the \textit{Soft Random Geometric Graph} (SRGG). In this model, nodes are connected at random according to a function of their spatial distance \cite{penrose2016connectivity}. 

Despite their numerous applications, the information theory of these random graphs is not well understood. In recent years there have been a growing number of works that have concentrated on this model or similar ones such as graphons. In \cite{janson2010graphons} a limit theorem for the entropy of $W$ random graphs (which include dense SRGGs as a special case) was proven. Since then, there have been a number of works attempting to bound or quantify the entropy, such as \cite{coon2018entropy, baker2026entropy, coon2018conditional, vippathalla2026entropy}. Also, the compression of random graphs has gained popularity recently. Models such as sparse Erd\H os-R\'enyi graphs \cite{choi_structural_entropy}, stochastic block models \cite{Abbe2016GraphClusters, martin_rate_distortion_sbm, bhatt2023universal} and preferential attachment graphs \cite{luczak_pag} have all been considered. Compression of the SRGG (or generalisations) has been studied in \cite{delgosha_universal_compression, delgosha2022universal, vippathalla2024lossy}. In works such as \cite{choi_structural_entropy, mihai2021structural}, the goal is to compress the \textit{structure} of the graph, whereas our aim is to compress the \textit{labelled} graph. Reference \cite{Paton2022labelled_entropy} gives a detailed treatment of the differences between these regimes. 

Crucially, all of these works focus on compression in the centralised setting, where the encoder sees the entire graph. However, in many applications, it is perhaps more realistic to assume that nodes, or blocks of nodes only have access to partial information about the network. In \cite{delgosha2021distributed}, distributed compression of Erd\H os-R\'enyi and configuration model graphs is studied, but distributed compression of the SRGG is as yet unstudied. This regime is of interest in SRGGs because the spatial nature of the graph can induce natural limits on how much of the graph can be observed. We consider the regime where we have a fixed, arbitrary number of encoders which can access different sections of the graph. Beyond graph compression, there are examples of situations where a central controller needs (at least partial) topological knowledge from distributed network locations. For example, sink nodes or central controllers in wireless sensor networks that make routing, scheduling, interference management or topology control decisions \cite{tunca2013distributed}.

Our other source of novelty is the specific regime in which we study the SRGG. In the above works, when the SRGG (or $W$ random graph) is made sparse (the expected degree of a node is sublinear in the total number of nodes), a fraction of the edges of a dense ensemble are removed randomly. However, this places less emphasis on the geometry, since the length scale of edges does not scale with the number of nodes. In our work, we study the more common regime, for example in the work of Penrose \cite{penrose2016connectivity}, where the `connection range' of each node is decreasing in the number of nodes. Equivalently, the size of the domain in which we embed the graph is growing as we add more nodes. We make this idea more precise later. To the best of our knowledge, this regime has not yet been studied in an information-theoretic context.

We also note that the SRGG is an example of a \textit{non-standard source}, in that it cannot be written as a string of stationary, ergodic or Markovian symbols. We must therefore employ the ideas of information spectrum theory to study its information theoretic properties. In this work, we contribute the following. \begin{enumerate}
    \item A limit theorem for the entropy rate of the SRGG,
    \item an Asymptotic Equipartition Property (AEP) for the SRGG,
    \item a complete characterisation of the Slepian-Wolf region for the distributed compression of the SRGG.
\end{enumerate} In Section \ref{sec:prelims} we formally introduce the model, the information spectrum, and the distributed compression setting. Section \ref{sec:results} presents our main results, and Section \ref{sec:proofs} contains proof ideas. The full proofs, and plots of the rate region may be found in the extended version on arXiv \cite{baker2026entropydistributedsourcecoding}. Finally, Section \ref{sec:conclusion} concludes.

\section{Model and Problem Formulation}
\label{sec:prelims}

\subsection{Soft Random Geometric Graphs}

\begin{definition}[Soft Random Geometric Graph]
    Let $Z_1,...,Z_n$ be independent, uniformly distributed points in a compact set $\mathcal{K} \subset \mathbb{R}^d$ with unit volume, a non-empty interior, whose boundary is piecewise smooth, of finite length and has a finite number of corners. The Soft Random Geometric Graph (SRGG) $G_n = (V_n, E_n)$ is the graph with vertex set $V_n = [n] := \{1,2,...,n\}$, and edge set formed by connecting each pair of vertices $i, j \in V_n$ independently with probability $p_n(R_{ij})$, where $R_{ij} := \|Z_i-Z_j\|$ is the Euclidean distance between nodes $i$ and $j$, and $p_n : [0,\infty) \to [0,1]$ is the `connection function'. Let $\graph_n(\mathcal{K})$ be the set of all SRGGs on $n$ nodes embedded in $\mathcal{K}$, then $\graph_n(\mathcal{K})$ is called the ensemble of SRGGs. We will denote this by $\graph_n$ in the remainder of this paper for ease of notation.
\end{definition}
In this work, we will consider specifically the following form of connection function. Let $\{s(n)\}_{n=1}^\infty$ be a positive, decreasing sequence of real numbers. We call $s(n)$ the \textit{sparsity} of the SRGG. Then, we focus on connection functions $p_n(r) = p(r/s(n))$, where $p : [0,\infty) \to [0,1]$. The connection function $p$ is usually non-increasing, with a common example being the Rayleigh fading connection function $p(r/s(n)) = \exp(-(r/s(n))^\eta)$, where $\eta > 0$ controls the decay of connectivity. This can be thought of as decreasing the `connection range' of each node as $n\tti$, that is, the length scale at which two nodes are likely to be connected is decreasing with $n$. We note that this is a different regime to the other commonly considered SRGG, where $p_n(r) = s(n)p(r)$. We will require that $1 \gg s(n) \gg \left(\frac{\log n}{n}\right)^{1/d}$. This requirement has two interpretations. The first is that this is the regime where the SRGG is connected with high probability as $n\tti$ \cite{penrose2016connectivity}. The second is that the entropy of the labelled graph is asymptotically equivalent to the entropy of the unlabelled graph.

\begin{definition}[Entropy of the SRGG]
    The Shannon Entropy of the SRGG is denoted by $H(\graph_n)$, and is given by
    \begin{equation}
        H(\graph_n) := - \sum_{G \in \graph_n} \prob{G_n}\log \prob{G_n}
    \end{equation}
    where the logarithm is to base 2, and we take the convention that $0 \log 0 = 0$. Also, if $Z_n$ is the random vector containing the locations of each node, then the conditional entropy $H(\graph_n|\mathcal{Z}_n) = \mathbb{E}[H(\graph_n | Z_n)]$ is given by
    \begin{equation}
        H(\graph_n|\mathcal{Z}_n) = \binom{n}{2}\int_0^\infty f_{\mathcal{K}}(r)h_2(p_n(r))dr
    \end{equation}
    where $f_{\mathcal{K}}$ is the probability density of pairwise distances in $\mathcal{K}$, and $h_2(p) = -p\log p-(1-p)\log(1-p)$ is the binary entropy function.
\end{definition}

We treat the following as a standing assumption for the remainder of the paper.

\begin{assumption}
    \label{ass:1}
    The connection function $p_n$ is given by $p_n(r) = p(r/s(n))$, where $p$ is a H\"older continuous function, that is, for every $(x,y),(x',y') \in \mathcal{K}^2$,
    \begin{equation}
        |p(x,y) - p(x',y')| \leq L\|(x,y)-(x',y')\|^\alpha
    \end{equation}
    with $L, \alpha > 0$ as constants. We assume that, if $d$ is the dimension of the space $\mathcal{K} \subset \mathbb{R}^d$,
        \begin{equation}
            \label{ass:var_1}
            \int_0^\infty r^{2d}h_2(p(r)) dr < \infty
        \end{equation}
        and
        \begin{equation}
            \label{ass:var_2}
            \int_0^\infty r^{d-1}p(r)(1-p(r))\log^2\left(\frac{p(r)}{1-p(r)}\right)dr <\infty
        \end{equation}
    We also define the following integral
    \begin{equation}
        \label{ass:h_star}
        h^* := \kappa_{d-1}\int_0^\infty r^{d-1}h_2(p(r))dr < \infty
    \end{equation}
    where $\kappa_{d-1}$ is the volume of the unit $d-1$ ball. 
\end{assumption}

\subsection{The Information Spectrum}

As briefly mentioned in the introduction, the SRGG is what is known as a `non-standard source'. In classical Information Theory, sources are considered to consist of a sequence of i.i.d. symbols, or at the very least, to be stationary and ergodic. However, the SRGG has neither of these properties. Instead, we must consider the graph $G_n$ in its entirety. \textit{Information spectrum theory} was developed to deal with general sources such as the SRGG. The theory is developed around \textit{spectral entropy rates}. We first define the \textit{limit superior in probability}, and the \textit{limit inferior in probability}. Let $(X_1,X_2,...)$ be a real-valued sequence of random variables, then
\begin{equation}
    \plimsup_{n\tti} X_n := \inf \left\{\alpha | \lim_{n\tti}\prob{X_n > \alpha} =0\right\}
\end{equation}
\begin{equation}
    \pliminf_{n\tti} X_n := \sup \left\{\beta | \lim_{n\tti}\prob{X_n < \beta} =0\right\}
\end{equation}

If $\pliminf_{n\tti} X_n = \plimsup_{n\tti} X_n = X^*$, we say that $\plim_{n\tti} X_n = X^*$, which is equivalent to $X_n \to X^*$ in probability as $n\tti$ \cite[\S 1.3]{han_book}.

\begin{definition}[Spectral Entropy Rates]
    Let $\mathbf{G} = \{G_n\}_{n=2}^\infty$ denote a general source (see \cite[\S 1.3]{han_book}), where $G_n$ is a $\graph_n$-valued random variable. We define the spectral sup-entropy rate of $\mathbf{G}$ as
    \begin{equation}
        \overline{H}(\mathbf{G}) = \plimsup_{n\tti} \frac{1}{\binom{n}{2}s(n)^d}\log \frac{1}{\prob{G_n}}
    \end{equation}
    Similarly, the spectral inf-entropy rate is defined as
    \begin{equation}
        \underline{H}(\mathbf{G}) = \pliminf_{n\tti} \frac{1}{\binom{n}{2}s(n)^d}\log \frac{1}{\prob{G_n}}
    \end{equation}
    If $\overline{H}(\mathbf{G}) = \underline{H}(\mathbf{G})$, then the limit in probability exists, and the spectral entropy rate
    \begin{equation}
        H(\mathbf{G}) = \plim_{n\tti} \frac{1}{\binom{n}{2}s(n)^d}\log \frac{1}{\mathbb{P}(G_n)}
    \end{equation}
    is well-defined.
\end{definition}

These quantities are similar to the spectral entropy rates used in information spectrum theory for general sources \cite{han_book}. The difference is in the normalisation. Normalising by the blocklength $\binom{n}{2}$, would give $\overline{H}(\mathbf{G})=0$. This is not a very interesting result, since the sparsity induced by $s(n)$ means the graph has a vanishing density, and so the majority of edges are absent in the $n\to\infty$ limit. Our definition may be thought of as an `effective' spectral entropy rate, capturing the leading order behaviour of compression length, rather than simply dividing through by the number of symbols in the source. It is simple to check that these quantities retain their operational meaning after replacing the standard scaling of $1/n$ by $1/(\binom{n}{2}s(n)^d)$. \\

\vspace{-0.2cm}
\section{Main Results}
\label{sec:results}

We first establish the following limit theorem.

\begin{theorem}[Limit Theorem]
    \label{thm:limit}
    Let $\graph_n$ be the ensemble of SRGGs with connection function $p_n(r) = p(r/s(n))$. Then,
    \begin{equation}
        \lim_{n\tti} \frac{H(\graph_n)}{\binom{n}{2}s(n)^d} = h^*
    \end{equation}
\end{theorem}

The next result strengthens the above to an AEP, showing that the spectral sup-entropy rate is equal to the spectral inf-entropy rate.

\begin{theorem}[Asymptotic Equipartition Property]
    \label{thm:aep}
    If $G_n$ is a random SRGG sampled from $\graph_n$,
    \begin{equation}
        \plim_{n\tti} \frac{1}{\binom{n}{2}s(n)^d} \log \frac{1}{\prob{G_n}} = h^*
    \end{equation} 
    In other words,
    \begin{equation}
        \overline{H}(\mathbf{G}) = \underline{H}(\mathbf{G}) = H(\mathbf{G}) = h^*
    \end{equation}
    This implies the existence of sets defined for all $\epsilon > 0$,
    \begin{equation}
        \mathcal{T}_n^\epsilon := \left\{G_n : \left|\log \frac{1}{\prob{G_n}} - H(\graph_n)\right| \leq \binom{n}{2}s(n)^d \epsilon\right\} 
    \end{equation}
    such that
    \begin{enumerate}
        \item $\prob{G_n \in \mathcal{T}_n^\epsilon} \to 1$ as $n\tti$
        \item For each $G_n \in \mathcal{T}_n^\epsilon$, $2^{-\binom{n}{2}(s(n)^dh^* +\epsilon)} \leq \prob{G_n} \leq 2^{-\binom{n}{2}s(n)^d(h^*- \epsilon)}$
        \item For every $\delta > 0$, there exists an $n^*$ so that for all $n > n^*$, we have
        \begin{equation}
            (1-\delta)2^{\binom{n}{2}s(n)^d(h^*-\epsilon)} \leq |\mathcal{T}_n^\epsilon| \leq 2^{\binom{n}{2}s(n)^d(h^*+\epsilon)}
        \end{equation}
    \end{enumerate}
\end{theorem}

We now introduce the distributed compression scheme for the SRGG. For each node $i \in [n]$, let $N_i$ be a $\{0,1\}^{(n-1)}$-valued random vector representing the neighbourhood of $i$ in $G_n$. That is, $N_i = (X_{ij})_{j \neq i}$, where $X_{ij} = 1$ if $i$ and $j$ share an edge, and is $0$ otherwise. In this paper we consider $X_{ij}$ to be the same random variable as $X_{ji}$. The distributed compression problem is formulated as follows. Let $L \in \mathbb{N}$ be a constant. We partition the graph $G_n$ into $L$ `blocks', given by
\begin{equation}
    B^{(n)}_{l} = \left\{N_i : (l-1)n/L < i \leq ln/L\right\}
\end{equation}
for $l = 1,...,L$. Denote by $\mathcal{B}_l^{(n)}$ the set of all realisations of $B^{(n)}_l$ for each $l$. The distributed compression problem asks whether there exists a sequence of encoders $\{(\phi_1^{(n)},...,\phi_L^{(n)})\}_{n=1}^\infty$, where $\phi_l^{(n)} : \mathcal{B}_l^{(n)} \to \mathcal{M}_l^{(n)}$, and a sequence of decoders $\{\psi^{(n)}\}_{n=1}^\infty$ with $\psi^{(n)}: \prod_{l=1}^n\mathcal{M}_l^{(n)} \to \graph_n$ such that
\begin{equation*}
    P_E^{(n)} := \prob{\psi^{(n)}(\phi_1^{(n)}(B^{(n)}_1),...,\phi_L^{(n)}(B^{(n)}_L)) \neq G_n} \to 0
\end{equation*}
as $n\tti$. We say an $L$-tuple of rates $(R_1,...,R_L) \in [0,\infty)^{L}$ is \textit{achievable} if 
\begin{equation*}
    \limsup_{n\tti} \frac{L}{n(n-1)s(n)^d} \log |\mathcal{M}_l^{(n)}| \leq R_l
\end{equation*}
for each $l=1,...,n$, and $P_E^{(n)} \to 0$ as $n\tti$. The above normalisation can be understood as $L/n$ for the size of the block, and $1/((n-1)s(n)^d)$ to normalise the per-neighbourhood entropy (consisting of $n-1$ symbols each). For this case, we have the following result
\begin{theorem}[Distributed Compression Rate Region]
    \label{thm:partial_dsc}
    In the distributed compression scheme as described above, a rate tuple $(R_1,...,R_L)$ is achievable if and only if
    \begin{equation}
        \sum_{l \in \Lambda} R_l \geq \frac{|\Lambda|^2}{L}\frac{h^*}{2} 
    \end{equation}
    for every $\Lambda \subseteq [L]$.
\end{theorem}
\begin{remark}
    In particular, when $\Lambda = [L]$, we require $\frac{1}{L}\sum_{l=1}^L R_l \geq \frac{h^*}{2}$. The factor of $1/2$ comes from the fact that the information about each edge is stored twice overall. To see this, note that the variable $X_{12}$ is in both $N_1$, and $N_2$. 
\end{remark}

\section{Proof Ideas}
\label{sec:proofs}

\subsection{Theorem \ref{thm:limit}}
We first prove that $H(\graph_n) = H(\graph_n|\mathcal{Z}_n) + O(n\log n)$. We follow the idea of \cite[Theorem D.5]{janson2010graphons}, which establishes this result in the dense regime ($s(n) = \Theta(1)$). It follows by discretising the domain $\mathcal{K}$ into $m^d \in \mathbb{N}$ boxes, and conditioning on the box that contains each node. The proof of our result requires $m$ to vary with $n$, which is not the case for the dense regime. Explicitly, if $M$ is the random vector collecting these labels, which has entropy $H(M) \leq nd\log m$,
\begin{equation}
    H(\graph_n|\mathcal{Z}_n) \leq H(\graph_n) \leq H(\graph_n|M) + H(M)
\end{equation}
\begin{align}
    \Rightarrow &H(\graph_n)-H(\graph_n|\mathcal{Z}_n) \leq nd\log m + H(\graph_n|M) - H(\graph_n|\mathcal{Z}_n) \nonumber \\
    &\leq nd\log m + \binom{n}{2}(H(X_{12}|M) - H(X_{12}|\mathcal{Z}_n))
\end{align}
Then, letting $m$ grow with $n$, the aim is to show that the second term is $o(s(n)^d)$. $H(X_{12}|M)$ is the entropy of an edge existing with respect to an `averaged' version of $p$, and so using the H\"older continuity of $p$ it can be shown that\vspace{2cm}
\begin{equation}
    H(X_{ij}|M) - H(X_{ij}|\mathcal{Z}_n) \leq (1+o(1)) C\frac{\log (s(n)m)}{s(n)^{\alpha}m^\alpha}.
\end{equation}
for some $C > 0$. This means that we can choose any $m = m(n)$ such that $\log(m) \ll ns(n)^d$ and $m s(n)^{(d+\alpha)/\alpha}\gg 1$ to get
\begin{equation}
    \frac{H(\graph_n)-H(\graph_n|\mathcal{Z}_n)}{\binom{n}{2}s(n)^d} \leq \frac{nd\log m}{\binom{n}{2}s(n)^d}+(1+o(1)) C\frac{\log (s(n)m)}{s(n)^{d+\alpha}m^\alpha}
\end{equation}
which goes to 0 as $n\tti$. To complete the result, we use Theorem 1 from \cite{baker2026entropy}. Under the assumption \eqref{ass:h_star}, this gives
\begin{equation}
    \frac{H(\graph_n|\mathcal{Z}_n)}{\binom{n}{2}s(n)^d} \overset{n\tti}{\to} h^*.
\end{equation}
Combining with the above argument yields the result. \qed

\subsection{Theorem \ref{thm:aep}}

We begin by writing
\begin{align}
    |-\log \prob{G_n} - H(\graph_n)| &\leq |H(\graph_n) - H(\graph_n|\mathcal{Z}_n)| \nonumber \\ &+ |-\log \prob{G_n} - H(\graph_n | \mathcal{Z}_n)| \nonumber 
\end{align}
By Theorem \ref{thm:limit} the first term is $o(n^2s(n)^d)$, and for the second term, we use Markov's inequality to give
\begin{align}
    \mathbb{P}(|-\log\space &\prob{G_n} - H(\graph_n|\mathcal{Z}_n)| > \binom{n}{2}s(n)^d\epsilon) \nonumber \\&\leq \frac{\mathbb{E}_{G_n}[|-\log \prob{G_n} - H(\graph_n|\mathcal{Z}_n)|]}{\binom{n}{2}s(n)^d\epsilon} \nonumber \\
    &= \frac{\mathbb{E}_{G_n,Z_n}[|-\log \prob{G_n} - H(\graph_n|\mathcal{Z}_n)| | Z_n]}{\binom{n}{2}s(n)^d\epsilon} \nonumber \\ 
    \label{eq:split_bound}
    &\leq \frac{\mathbb{E}_{G_n,Z_n}[|-\log \prob{G_n|Z_n} - H(\graph_n|\mathcal{Z}_n)|]}{\binom{n}{2}s(n)^d\epsilon} \nonumber \\ &+ \frac{\mathbb{E}_{G_n,Z_n}[|\log \prob{G_n|Z_n} - \log \prob{G_n}|]}{\binom{n}{2}s(n)^d\epsilon}
\end{align}

To show that the first term goes to zero, we need the following lemma.
\begin{lemma}[Conditional AEP]
    \label{thm:conditional_aep}
    Let $\{(G_n,Z_n)\}_{n=2}^\infty$ be a sequence of pairs of random sparse SRGGs and the node locations that generated it. Then,
    \begin{equation}
        s(n)^{-d}\binom{n}{2}^{-1}\left(-\log \prob{G_n|Z_n} - H(\graph_n|\mathcal{Z}_n)\right) \to 0
    \end{equation}
    in expectation as $n\tti$.
\end{lemma}

 This is easily proved using the weak law of large numbers. The only condition we need to verify is that:
\begin{align*}
    &\frac{\Var(\log \prob{G_n|Z_n})}{\binom{n}{2}^2s(n)^{2d}} \nonumber \\ &= \frac{\Var(X_{12}\log \prob{X_{12}|Z_n}+ (1-X_{ij})\log(1-\mathbb{P}(X_{12}|Z_n)))}{\binom{n}{2}s(n)^{2d}}
\end{align*}
is finite, which allows us to use Chebyshev's inequality. This is proven by a similar argument to the proof of \cite[Theorem 1]{baker2026entropy} cited above, and uses assumptions \eqref{ass:var_1} and \eqref{ass:var_2}. This condition also allows us to upgrade convergence in probability (given by WLLN) to convergence in expectation via de la Vall\'ee Poussin's criterion for uniform integrability.

The second term goes to 0 via the following argument. Let
\begin{equation}
    r_n = r(G_n, Z_n) = \frac{\prob{G_n|Z_n}}{\prob{G_n}}.
\end{equation}
Then, we have
\begin{equation*}
    \mathbb{E}_{G_n,Z_n}[|\log \prob{G_n|Z_n} - \log \prob{G_n}|] = \mathbb{E}_{G_n,Z_n}[|\log r_n|]
\end{equation*}
\begin{equation}
    = \mathbb{E}_{Z_n}\left[\sum_{G_n \in \graph_n} \prob{G_n|Z_n} |\log r_n|\right]
\end{equation}
\begin{equation}
    = \mathbb{E}_{Z_n}\left[\sum_{G_n \in \graph_n} \prob{G_n} |r_n \log r_n| \right]
\end{equation}
Then, we make the following two observations. First, $|x| = x - 2x\mathds{1}(x < 0)$, and second, for $x \in (0,1)$, $0 < -x\log x < \frac{1}{e}\log e$. Therefore,
\begin{align}
    \mathbb{E}_{G_n,Z_n}[|\log r_n|] \leq \mathbb{E}_{Z_n}\left[\sum_{G_n \in \graph_n} \prob{G_n} r_n \log r_n \right] \nonumber \\ - 2\mathbb{E}_{Z_n}\left[\sum_{G_n \in \graph_n} \prob{G_n} r_n \log r_n \Bigg| r_n < 1 \right] \nonumber \\
    \leq \mathbb{E}_{Z_n}\left[\sum_{G_n \in \graph_n} \prob{G_n} r_n \log r_n \right] \nonumber \\+ \frac{2}{e}\log(e) \mathbb{E}_Z\left[\sum_{G_n \in \graph_n} \prob{G_n} \right] \nonumber \\
    = \mathbb{E}_{G_n,Z_n}\left[\log \frac{\prob{G_n|Z_n}}{\prob{G_n}}\right] + \frac{2}{e}\log e
\end{align}
\begin{equation}
    = H(\graph_n) - H(\graph_n|\mathcal{Z}_n) + \frac{2}{e}\log e = o(n^2s(n)^d)
\end{equation}
This places an $o(1)$ bound on the probability that $|\log(1/\prob{G_n}) - H(\graph_n)| > \binom{n}{2}s(n)^d\epsilon$ for any $\epsilon>0$. Combining this with the fact that $(\binom{n}{2}s(n)^d)^{-1}H(\graph_n) \to h^*$ as $n\tti$, we get $\plim_{n\tti} \frac{1}{\binom{n}{2}s(n)^d}\log \frac{1}{\prob{G_n}} = h^*$. The properties of the typical set are derived using simple standard techniques.

\subsection{Theorem \ref{thm:partial_dsc}}

\begin{lemma}
    \label{thm:nbhd_reduction}
    Let $S_n \subseteq V_n$, and $S_n^c := V_n \setminus S_n$ then
    \begin{equation}
        H(N_{S_n}|N_{S_n^c}) = H(\graph_n[S_n])
    \end{equation}
    where $\graph_n[S_n]$ is the ensemble of subgraphs $G_n[S_n]$ of $G_n$ restricted to node set $S$. This implies that if $\{S_n\}_{n=1}^\infty$ is a sequence of subsets of $V_n$ indexed by $n$, such that $|S_n| \to \infty$ as $n\tti$, then
    \begin{equation}
        \label{eq:nbhd_aep}
        \plim_{n\tti} \frac{1}{\binom{|S_n|}{2}s(n)^d}\log \frac{1}{\prob{N_{S_n}|N_{S_n^c}}} = h^* 
    \end{equation}
\end{lemma}

\begin{proof}
    For the first part, fix $n$, and let $k < n$. We will show the first part for ${S} = \{1,..,k\}$,  which generalises to any set ${S}$ by re-labelling the graph. Writing out $N_{{S}}$ in full, we have, recalling that $X_{ij}$ is identified with $X_{ji}$,
    \begin{equation}
        N_{{S}} = \{X_{ij}\}_{\substack{i \in {S}, 1 \leq j \leq n\\
                                    i \neq j}}, N_{{S}^c} = \{X_{ij}\}_{\substack{i \notin {S}, 1 \leq j \leq n\\
                                    i \neq j}}
    \end{equation}
    Then, the entropy $H(N_{S}|N_{S^c
    })$ is the entropy of $A_{{S}} := N_{S} \setminus (N_{S} \cap N_{{S}^c})$ conditioned on $\widetilde{A}_{S} := N_{{S}^c} \setminus (N_{{S}} \cap N_{{S}^c})$, since the edges in the intersection do not contribute to the entropy. The key observation is that if $X_{ij} \in A_{S}$, then there does not exist an edge $X_{uv}$ in $\widetilde{A}_{S}$ with $u = i$ or $v = j$. Therefore, every edge variable in $A_{S}$ is independent of all those in $\widetilde{A}_{S}$. Then,
    \begin{equation*}
        H(N_S|N_{S^c}) = H(A_{S}|\widetilde{A}_{S}) = H(\{X_{ij}\}_{1\leq i<j \leq k}) = H(\graph_n[S])
    \end{equation*}
    as required. For the second part of the theorem, let $\{S_n\}$ be a sequence of sets with size proportional to $n$. Since each edge is added to the graph independently, the subgraph $G_n[S]$ is marginally distributed as a SRGG on $|S|$ nodes with sparsity $s(n)$ inherited from the parent graph $G_n$. Then, we can apply Theorem \ref{thm:aep} to get the result.
\end{proof}
\begin{proof}[Proof of Theorem \ref{thm:partial_dsc}]
    For $\Lambda \subseteq [L]$, let $B_{\Lambda}^{(n)} = \bigcup_{l\in \Lambda} B_l^{(n)}$, the union of the neighbourhoods of all the nodes in block $l$ for each $l \in \Lambda$. After establishing the above Lemma \ref{thm:nbhd_reduction}, for each $\Lambda \subseteq [L]$, and for all $\epsilon > 0$ we have a typical set $\mathcal{T}_n^\epsilon(\Lambda)$, which we define as
    \begin{align}
        &\mathcal{T}_n^\epsilon(\Lambda) := \Big\{(B_\Lambda^{(n)},B_{\Lambda^c}^{(n)}) :\nonumber \\ &\Big| \frac{1}{{\binom{{|\Lambda|n}/{L}}{2} s\left(n\right)^d}}\log {\mathbb{P}(B_{\Lambda}^{(n)}|B_{\Lambda^c}^{(n)})} 
        -h^* \Big| < \epsilon\Big\}
    \end{align}
    whose existence is guaranteed by the second part of Lemma \ref{thm:nbhd_reduction}.
    With this, we can apply the standard random binning argument, which we sketch here. Let $L$ be a fixed integer, then we design the codebook as follows for an $L$-tuple of rates $(R_1,...,R_L)$. For each $l = 1,...,L$, we generate assignments by independently assigning each $b_l^{(n)} \in \mathcal{B}_l^{(n)}$ a uniformly random element, called its index $\phi_l^{(n)}(b_l^{(n)})$, from the codebook
    \begin{equation*}
        \mathcal{M}_l^{(n)} := \{1,2,3,...,2^{\frac{n}{L}(n-1)s(n)^dR_l}\}.
    \end{equation*}
    Note that more than one realisation $b_l^{(n)}$ may be assigned the same index. We reveal these assignments to both the encoder, and the decoder. To encode a graph, each encoder $l$ sees its block, $b_l^{(n)}$. Then, the independent encodings are sent as the tuple $(\phi_1^{(n)}(b_1^{(n)}),...,\phi_L^{(n)}(b_L^{(n)}))$. The decoder receives this tuple, and searches for their own tuple $(\hat{b}_1^{(n)},...,\hat{b}_L^{(n)})$ (or equivalently a graph $\hat{G}_n$), such that $(\hat{b}_1^{(n)},...,\hat{b}_L^{(n)})$ is unique, typical, and $\phi_l^{(n)}(\hat{b}_l^{(n)}) = \phi_l^{(n)}(b_l^{(n)})$ for every $l$. If they cannot find such a tuple, an error is declared. Therefore, an error occurs if, first, the graph is atypical, or second, there exists an alternative decoding that is typical, and matches the encoding of the target graph. By Theorem \ref{thm:aep} the first error probability goes to 0. The probability of the second error event is bounded above by the sum over such $\Lambda$ of
    \begin{align}
        \prob{E_2(\Lambda)} = \mathbb{P}\Big(\exists &\hat{b}_\Lambda^{(n)}\neq B_\Lambda^{(n)} \space \text{ s.t. } (B_{\Lambda}^{(n)}, B_{\Lambda^c}^{(n)}) \in \mathcal{T}_n^\epsilon(\Lambda) \nonumber \\ &\text{ and } \phi_l^{(n)}(\hat{b}_l^{(n)}) = \phi_l^{(n)}(B_l^{(n)})\space \forall l \in \Lambda\Big) \nonumber 
    \end{align}
    Following the standard argument, this is bounded above by
    \begin{equation}
        \prob{E_2(\Lambda)} \leq 2^{-\frac{n(n-1)}{L}s(n)^d\sum_{l\in \Lambda}R_l}|\mathcal{T}_n^\epsilon(\Lambda)|
    \end{equation}
    \begin{equation}
        \leq 2^{-\frac{n(n-1)}{L}s(n)^d\sum_{l\in \Lambda}R_l + \binom{v}{2}s(n)^d(h^*+\epsilon)}
    \end{equation}
    where $v = |B_{\Lambda}^{(n)}| = |\Lambda|n/L$. From this we can see that a sufficient condition for achievability w.h.p. as $n\tti$ is
    \begin{equation}
        \sum_{l\in \Lambda} R_l \geq \frac{|\Lambda|^2}{L}\frac{h^*}{2}
    \end{equation}
    The proof of the converse is standard and may be found in the extended version. The important point is that
    \begin{align}
        &\lim_{n\tti} P_E^{(n)} \nonumber \\ &\geq \lim_{n\tti} \mathbb{P}\Big(\frac{1}{{\frac{n}{L}}(n-1)s(n)^d} \log \frac{1}{\mathbb{P}(B_{\Lambda}^{(n)}|B_{\Lambda^c}^{(n)})} \geq \sum_{l \in \Lambda} R_l + \gamma\Big) \nonumber \\
        &= \mathds{1}\left(\lim_{n\tti}\frac{2}{|\Lambda|}\frac{(n-1)}{\left(|\Lambda|n/{L}-1\right)}\left(\sum_{l \in S} R_l + \gamma\right) < h^*\right)
    \end{align}
    which follows from \cite[Lemma 7.2.2]{han_book}, the definition of $\plim$ and Lemma \ref{thm:nbhd_reduction}, and gives that any achievable rate tuple must satisfy
    \begin{equation}
        \sum_{l \in \Lambda} R_l \geq \frac{|\Lambda|^2}{L}\frac{h^*}{2}
    \end{equation}
    for the matching converse.
\end{proof}

\section{Conclusion}
\label{sec:conclusion}

In this paper, we have studied the information theoretic properties of the soft random geometric graph in the sparse super-connectivity regime. We gave limit results which characterise the entropy rate of the SRGG, and establish that it is an information-stable source, in that its information content converges to the entropy rate with high probability. Next, we found the Slepian-Wolf rate region for distributed compression of the SRGG where there are a finite number of blocks of nodes. For future research, a tractable distributed compression algorithm would be a practically useful direction. From a theoretical standpoint, it would be interesting to extend our results to the regime where the number of compression units scales with $n$, and compare this region to the finite block size we considered. Extensions to lossy distributed compression may be of interest, including Berger-Tung bounds for this problem for a range of distortion functions.


\section*{Acknowledgment}
 The authors would like to thank Yihan Zhang for his helpful discussion. OB was funded by the EPSRC Centre for Doctoral Training in Computational Statistics and Data Science (COMPASS), Grant Number EP/S023569/1.
 
\bibliographystyle{IEEEtran}
\bibliography{bibliofile}

\end{document}
